% MT17091C
% Mechanische Atemwegsverlegung
% 14.0
% 2013-11-30

:ref:`m-atemwegsverlegung`
Taktik 
\begin{inparadesc}
  \item[Schwer:] \Ca{Vitale Bedrohung!} Zeitkritisch; wenn
möglich rasche Entfernung des Fremdkörpers. 
\item[Mild:] Symptomatische
Therapie.
\end{inparadesc}


\begin{itemize}
  -   Atemwege soweit möglich freimachen
  -   O₂-Gabe hochdosiert (\VonBis{10}{15} L / min)
  -   Lagerung: Situationsgerecht, z.\,B. Oberkörper erhöht
  -   |TxBeiVit| |TxMassVitMK|
  -   Weitere Maßnahmen abhängig von der Beurteilung der
  Schwere der Atemwegsverlegung:
\end{itemize}

{\FontSansNarrowFamily
\noindent\begin{tabularx}{\linewidth}{XXX}
\toprule\addlinespace
\multicolumn{3}{c}{\TabSub{Beurteilung der Schwere der Atemwegsverlegung}}
\\\addlinespace\midrule\addlinespace
\TabSub{Schwer}

Kann nicht sprechen

Hustet ineffektiv
&
&
\TabSub{Mild}

Kann sprechen

Hustet effektiv
\\\addlinespace\cmidrule(lr){1-2}\cmidrule(lr){3-3}\addlinespace
\TabSub{Ohne Bewusstsein}
&
\TabSub{Bei Bewusstsein}
&
\TabSub{Bei Bewusstsein}
\\\addlinespace\cmidrule(lr){1-1}\cmidrule(lr){2-2}\cmidrule(lr){3-3}\addlinespace
Reanimation
&
5 Schläge zwischen Schulterblätter

5 Heimlich-Manöver (\textgreater{}\,1\,\MetricUnit{a}) oder Thoraxkompressionen
(\textless{}\,1\,\MetricUnit{a})

Wiederholung 
&
Husten lassen

Überwachung

Keine weitere Manipulation
\\\addlinespace\bottomrule\addlinespace
\end{tabularx}
}

Bei Schlägen
    zwischen die Schulterblätter soll der Patient
    nach Möglichkeit abwärts schauen (besonders relevant
    bei Kindern).

Durch Anwendung des Heimlich-Handgriffes kann es zu
schweren inneren Verletzungen kommen. Ein Patient ist
nach Anwendung unbedingt zu hospitalisieren!

\citep{Erc2005DeSec2}